\ncertchapter{17}{Kepler's Laws and Orbits}\chlabel{ch:kepler}


Johannes Kepler spent the first decades of the seventeenth century fitting Tycho Brahe's planetary
observations, the most precise of their age, and emerged with three laws. Planets move in ellipses with the
Sun at one focus. The line from the Sun to a planet sweeps out equal areas in equal times. The square of the
period is proportional to the cube of the orbit's size. The laws described the solar system almost
perfectly, but they were descriptions: Kepler could not derive them from any law of force, and so could
not say why the orbits were ellipses, or why the periods followed that particular power.

Newton explained all three, and this chapter is that explanation. The second law turns out to be nothing
but the conservation of angular momentum, which holds because gravity is a central force and exerts no
torque about the Sun. The third law follows in two lines from setting the gravitational force equal to the
centripetal force for a circular orbit. The first law, that orbits are ellipses, is the hardest of the
three. It is usually quoted without proof at this level, but there is a proof, due to Hamilton, that needs
nothing beyond rotating unit vectors (as in \chref{ch:circkin}) and the conservation of angular momentum;
it is given in \S\ref{sec:ellipse} for the reader who wants to see the whole chain closed.

Almost every tool the chapter uses was built earlier (angular momentum in \chref{ch:rotation},
centripetal force in \chref{ch:circdyn}, potential energy in \chref{ch:gravitation}), so the chapter
doubles as a review of the volume.

\section{The Three Laws}

Kepler's laws are facts extracted from observation. They are stated first, with the vocabulary of the
ellipse they need, and then each is derived from Newton's law.

\begin{keybox}{Kepler's laws}
\textbf{First (orbits):} every planet moves in an ellipse with the Sun at one focus.
\textbf{Second (areas):} the line from Sun to planet sweeps equal areas in equal times.
\textbf{Third (periods):} $T^2\propto a^3$, with $T$ the time for one orbit and $a$ the semi-major axis
(for a circle, $a=r$).
\end{keybox}

\begin{defbox}{The ellipse and its vocabulary}
An \emph{ellipse} is the set of points whose distances $r_1$ and $r_2$ from two fixed points, the
\emph{foci}, have a constant sum: $r_1+r_2=2a$. Half the longest diameter is the \emph{semi-major axis}
$a$; half the shortest is the \emph{semi-minor axis} $b$. Each focus lies a distance $ae$ from the centre,
where the \emph{eccentricity} $e$ (a pure number, $0\le e<1$; not the coefficient of restitution of
\chref{ch:collisions}) measures how stretched the ellipse is, and
$b=a\sqrt{1-e^2}$. For $e=0$ the foci merge and the ellipse is a circle of radius $a$. With the Sun at one
focus, the nearest point of the orbit is the \emph{perihelion}, at distance $r_p=a(1-e)$, and the farthest
is the \emph{aphelion}, at $r_a=a(1+e)$ (for a satellite of the Earth, \emph{perigee} and
\emph{apogee}). Adding these, $r_p+r_a=2a$: the semi-major axis is the average of the nearest and
farthest distances (Figure~\ref{fig:orbitellipse}). The other focus is empty.
\end{defbox}

\begin{figure}[htbp]
\centering
\tikzfig{orbit_ellipse}
\figcaption{fig:orbitellipse}{The vocabulary of an orbit, drawn for $e=0.6$. The Sun sits at a focus,
not at the centre, so the planet's distance varies from $r_p$ to $r_a$; the semi-major axis $a$ is their
average. Real planetary orbits are far rounder than this: the Earth's has $e=0.017$, Mercury's $0.21$,
while Halley's comet, with $e=0.97$, swings from inside Venus's orbit to beyond Neptune's.}
\end{figure}

An ellipse belongs to the family of \emph{conic sections}, the curves cut from a cone by a plane. Tilting
the plane steeper gives, in turn, the circle ($e=0$), the ellipse ($0<e<1$), the parabola ($e=1$) and the
hyperbola ($e>1$); the last two are open curves that never close. \S\ref{sec:ellipse} shows that a body
moving under gravity follows one of these four, which one being decided by its energy.

\subsection{The second law is angular-momentum conservation}

Kepler's second law, the law of equal areas, is the easiest of the three to explain, and the explanation needs
only one fact about gravity: that it points along the line joining the two bodies, which makes it a
\emph{central} force (\chref{ch:gravitation}). Recall from \chref{ch:rotation} the torque
$\vtau=\vr\ex\vF$ and the angular momentum $\vL=\vr\ex m\vv$ of a particle at position $\vr$ (here
measured from the Sun) with velocity $\vv$; both are vectors, and $\vtau=d\vL/dt$.

\begin{derivbox}{Kepler's second law in three lines}
Kepler found the area law empirically and thought it a fact about planets. It is nothing but conservation
of angular momentum, and it holds for \emph{any} central force.
\stepn{1} Gravity is central: $\vF$ points along $\vr$, straight at the Sun. So the torque about the Sun
is $\vtau=\vr\ex\vF=0$ (the cross product of parallel vectors vanishes).
\stepn{2} Zero torque means constant angular momentum: $\vL=\text{constant}$ (\chref{ch:rotation}).
Because $\vL=\vr\ex m\vv$ is perpendicular to both $\vr$ and $\vv$, a fixed $\vL$ also keeps the orbit
in one plane.
\stepn{3} In a small time $dt$ the planet moves by $d\vr=\vv\,dt$, sweeping the thin triangle with sides $\vr$
and $d\vr$; its area is half that of their parallelogram (\chref{ch:vectors}), $dA=\tfrac12|\vr\ex d\vr|$, so
\[
\frac{dA}{dt}=\tfrac12|\vr\ex\vv|=\frac{|\vr\ex m\vv|}{2m}=\frac{L}{2m}=\text{constant}.
\]
Equal areas in equal times \emph{is} $L=\text{constant}$, wearing geometry (Figure~\ref{fig:keplerareas}).
\end{derivbox}

\Needspace*{23\baselineskip}
\begin{figure}[htbp]
\centering
\tikzfig{kepler_areas}
\figcaption{fig:keplerareas}{The area law is angular-momentum conservation. A planet races through
perihelion (closest, fast) and dawdles at aphelion (farthest, slow) for the same reason a skater speeds
up pulling her arms in: $L=mvr_\perp$ is fixed, with $r_\perp$ the perpendicular distance from the Sun
to the line of the velocity, so small $r$ demands large $v$. The first and third laws
are special to $1/r^2$; the second holds for any central force.}
\end{figure}

\section{Circular Orbits}

The second law holds for any central force and says nothing about the inverse square. To reach the
third law, which does depend on it, start with the simplest orbit, a circle (an ellipse with $e=0$),
for which the speed is constant and the acceleration is purely centripetal. It is enough to derive
Kepler's third law and the orbital energy.

\begin{derivbox}{Orbital speed, period, and Kepler's third law}
Let a body of mass $m$ circle a much heavier one of mass $M$ (the Sun for a planet, the Earth for a
satellite) at radius $r$ from its centre. Gravity supplies exactly the centripetal force (\chref{ch:circdyn}):
\[
\frac{GMm}{r^2}=\frac{mv^2}{r}.
\]
\stepn{1} \textbf{Orbital speed.} Cancel one $r$ and an $m$: $v^2=\dfrac{GM}{r}$, so the orbital speed
(met for $r=R$ in \S\ref{sec:escape}) is $v_o=\sqrt{\dfrac{GM}{r}}$, \emph{slower} in higher orbits. The satellite's mass $m$ has cancelled: a
heavy satellite and a light one at the same radius need the same speed, which is why a spanner let go
by an astronaut stays beside the station instead of drifting down.
\stepn{2} \textbf{Period and Kepler's third law.} The period $T$ is the time for one orbit, and the
angular speed is $\omega=2\pi/T$ (\chref{ch:circkin}). Write the centripetal force as
$m\omega^2 r$: $\dfrac{GMm}{r^2}=m\dfrac{4\pi^2}{T^2}r$, so
\[
T^2=\frac{4\pi^2}{GM}\,r^3.
\]
$T^2\propto r^3$ exactly as Kepler measured, but now with the constant $4\pi^2/GM$ \emph{predicted},
not fitted. For an
orbit just above the surface, $r=R$ and $GM=gR^2$ give $T=2\pi\sqrt{R/g}\approx\SI{5080}{s}$, about
$85$ minutes, the shortest period any Earth satellite can have.
\stepn{3} \textbf{Where the 3 comes from.} It is $1+2$: the $2$ from the inverse-square law, the $1$ from
the circle's geometry. Change gravity to $1/r^3$ and the exponent changes, so the exponent is a
\emph{fingerprint} of the force law, and finding the same exponent from Mercury out to Neptune is one
way we know that gravity is inverse-square to high precision.
\stepn{4} \textbf{Two refinements, stated without proof.} For an elliptical orbit the same law holds with
$r$ replaced by the semi-major axis $a$: $T^2=4\pi^2a^3/GM$. And strictly the Sun is not fixed: both bodies
circle their common centre of mass (\chref{ch:com}), and $GM$ becomes $G(M+m)$. For the Earth this changes
$T$ by under two parts in a million; for Jupiter, by about one part in two thousand.
\end{derivbox}

The speed fixes the energy too, a result already met in the satellite example of \chref{ch:gravitation};
here we read what it means.

\begin{derivbox}{Orbital energy, and the paradox}
With $v^2=GM/r$ the kinetic energy is $K=\tfrac12 mv^2=\dfrac{GMm}{2r}$, and $U=-\dfrac{GMm}{r}$. Add them:
\[
E=K+U=\frac{GMm}{2r}-\frac{GMm}{r}=-\frac{GMm}{2r}=-K.
\]
The total energy is \emph{negative} (bound) and equal to \emph{minus} the kinetic energy. Read the
consequence carefully: a satellite in a \emph{lower} orbit moves \emph{faster} (larger $K$) but has
\emph{less} total energy (more negative $E$). So to climb to a higher orbit you must \emph{add} energy,
and the satellite ends up moving \emph{slower} (Figure~\ref{fig:orbitenergy}). Raising $E$ to zero
frees the satellite, so its binding energy is $-E=GMm/2r$, and the escape condition gives
$v_e=\sqrt2\,v_o$ at every radius, not only at the surface.
\end{derivbox}

\begin{figure}[htbp]
\centering
\tikzfig{orbit_energy}
\figcaption{fig:orbitenergy}{The orbital paradox: a higher orbit has more total energy yet less speed.
The lower panel shows why. For every circular orbit $K=-E$ and $U=2E$, so as $r$ grows and $E$ rises
towards zero, $K$ must fall. This is why spacecraft rendezvous is counter-intuitive: to catch something
ahead of you, you \emph{slow down}, dropping into a lower and faster orbit, and overtake from below.
For the same reason, air drag on a low satellite, which removes energy, makes it sink and speed up.}
\end{figure}

\why{A satellite is a projectile (\chref{ch:projectiles}) and nothing more: it was thrown, and it is
falling towards the Earth without pause. It never lands because it moves sideways fast enough that the
Earth curves away beneath it as fast as it falls. As \chref{ch:gravitation} found for the ISS, gravity is
still $89\%$ of its surface value there; the astronauts float because they, the station and the drifting
pen all fall together.}

\begin{exbox}{The geostationary satellite}
To hang over one spot, the period must equal one rotation of the Earth, a \emph{sidereal} day,
$T=\SI{86164}{s}$ (not 24 hours, because the Earth also advances along its orbit). Using the measured
$GM=\SI{3.99e14}{m^3/s^2}$ and $R=\SI{6370}{km}$ (the rounded $gR^2$ of earlier examples would put $r$
about \SI{100}{km} too high):
\[
r^3=\frac{GM\,T^2}{4\pi^2}=\SI{7.50e22}{m^3}\ \Rightarrow\ r=\SI{42200}{km},\quad h=r-R=\SI{35800}{km}.
\]
About $5.6$ Earth-radii up, a tenth of the way to the Moon: a single ring, shared by every
geostationary communications satellite. Its speed is
$v=\sqrt{GM/r}=\SI{3.1}{km/s}$, against the ISS's \SI{7.7}{km/s}: the higher satellite moves less than
half as fast. \ans{\SI{35800}{km}; \SI{3.1}{km/s}}
\end{exbox}

\section{Elliptical Orbits: Speeds and Energy}\label{sec:ellipticorbits}

Most real orbits are not circles, and in an ellipse neither the speed nor the distance is constant, so
setting gravity equal to $mv^2/r$ no longer works. Two conservation laws take its place: angular
momentum (because the force is central) and energy (because gravity is conservative). Together they
give the speeds at the two ends of the orbit, and they show that the energy depends only on $a$.

\begin{derivbox}{Speeds at perihelion and aphelion, and $E=-GMm/2a$}
\stepn{1} \textbf{Angular momentum.} At perihelion and aphelion the velocity is perpendicular to $\vr$,
so $L=mvr$ there exactly, and conservation of $L$ gives $v_pr_p=v_ar_a$.
\stepn{2} \textbf{Energy.} $\tfrac12mv_p^2-\dfrac{GMm}{r_p}=\tfrac12mv_a^2-\dfrac{GMm}{r_a}$. Multiply by $2/m$
and substitute $v_a=v_pr_p/r_a$:
\[
v_p^2\,\frac{r_a^2-r_p^2}{r_a^2}=2GM\,\frac{r_a-r_p}{r_pr_a}
\quad\Longrightarrow\quad
v_p^2=\frac{2GM\,r_a}{r_p(r_p+r_a)},
\]
after cancelling $r_a-r_p$ (with $r_a^2-r_p^2=(r_a-r_p)(r_a+r_p)$).
\stepn{3} \textbf{The total energy.} Put this back into $E=\tfrac12mv_p^2-GMm/r_p$:
\[
E=GMm\left[\frac{r_a}{r_p(r_p+r_a)}-\frac{1}{r_p}\right]=GMm\,\frac{r_a-(r_p+r_a)}{r_p(r_p+r_a)}=-\frac{GMm}{r_p+r_a}=-\frac{GMm}{2a},
\]
using $r_p+r_a=2a$. This is the circular result $-GMm/2r$ with $r$ replaced by the semi-major axis: orbits with the same $a$
have the same energy (and, by Kepler's third law, the same period), however stretched they are.
\end{derivbox}

\begin{exbox}{A satellite on an elongated orbit}
A satellite's perigee is \SI{7000}{km} and its apogee \SI{21000}{km}, both measured from the Earth's
centre. Find its speeds at the two points and its period ($GM=\SI{3.99e14}{m^3/s^2}$).

\stepn{1} \textbf{Perigee speed} from the derivation:
$v_p=\sqrt{\dfrac{2(3.99\times10^{14})(2.1\times10^{7})}{(7.0\times10^{6})(2.8\times10^{7})}}=\SI{9.25}{km/s}$.
\stepn{2} \textbf{Apogee speed} from angular momentum: $v_a=v_p\,r_p/r_a=9.25/3=\SI{3.08}{km/s}$.
\stepn{3} \textbf{Period.} $a=\tfrac12(7000+21000)=\SI{14000}{km}$, so
$T=2\pi\sqrt{a^3/GM}=\SI{1.65e4}{s}\approx\SI{4.6}{h}$.

Compare the circular speed at \SI{7000}{km}, $\sqrt{GM/r_p}=\SI{7.55}{km/s}$: going faster than that at
the same point stretches the circle into this ellipse, which is exactly what Newton's cannon
(\chref{ch:gravitation}) showed. \ans{\SI{9.25}{km/s}; \SI{3.08}{km/s}; \SI{4.6}{h}}
\end{exbox}

\section{Why the Orbits Are Ellipses}\label{sec:ellipse}

Everything so far has assumed Kepler's first law: the orbits have been taken to be ellipses, and the
last section used only their ends. The remaining question is why an inverse-square force produces
exactly this curve and no other.

\scopebox{This section is beyond the examination syllabus, which asks only that the first law be stated.
It is included because it completes the argument of the chapter: all three of Kepler's laws then follow
from Newton's second law with the force $\vF=-\dfrac{GMm}{r^2}\uhat{r}$, and nothing else.}

\begin{derivbox}{Kepler's first law, by Hamilton's method}
Put the Sun at the origin and describe the planet by polar coordinates $(r,\theta)$ in the plane of its
orbit, with the unit vectors $\uhat{r}=(\cos\theta,\sin\theta)$ and $\uhat{\theta}=(-\sin\theta,\cos\theta)$.
Differentiating these with respect to $\theta$ gives $\dfrac{d\uhat{\theta}}{d\theta}=-\uhat{r}$.
\stepn{1} \textbf{Newton's second law} reads $m\dfrac{d\vv}{dt}=-\dfrac{GMm}{r^2}\,\uhat{r}$.
\stepn{2} \textbf{Trade time for angle.} Only the sideways velocity component, $r\,d\theta/dt$, contributes to
$|\vr\ex m\vv|$, so the angular momentum is $L=mr\cdot r\dfrac{d\theta}{dt}=mr^2\dfrac{d\theta}{dt}$, a
constant; hence
$dt=\dfrac{mr^2}{L}\,d\theta$, and
\[
\frac{d\vv}{d\theta}=\frac{d\vv}{dt}\,\frac{dt}{d\theta}=-\frac{GM}{r^2}\,\uhat{r}\,\frac{mr^2}{L}
=-\frac{GMm}{L}\,\uhat{r}=\frac{GMm}{L}\,\frac{d\uhat{\theta}}{d\theta}.
\]
The $r^2$ has cancelled: that is the inverse-square law doing the work.
\stepn{3} \textbf{Integrate.} Both sides are derivatives with respect to $\theta$, so
\[
\vv=\frac{GMm}{L}\,\uhat{\theta}+\vc{C},
\]
with $\vc{C}$ a constant vector. Choose the axes so that $\vc{C}$ points along $+y$, with magnitude $C$.
\stepn{4} \textbf{Take the $\uhat{\theta}$ component.} On the left, the sideways part of the velocity is
$r\,d\theta/dt=L/(mr)$, from step~2. On the right, $\uhat{\theta}\cdot\uhat{\theta}=1$ and
$\vc{C}\cdot\uhat{\theta}=C\cos\theta$. So
\[
\frac{L}{mr}=\frac{GMm}{L}+C\cos\theta
\quad\Longrightarrow\quad
r=\frac{p}{1+e\cos\theta},\qquad p=\frac{L^2}{GMm^2},\quad e=\frac{CL}{GMm}.
\]
Here $p$ is a length, the \emph{semi-latus rectum} (not a momentum), and $e$ a pure number; step~5 shows
that $e$ is the eccentricity defined earlier in this chapter.
\stepn{5} \textbf{Recognise the curve.} Write $r+e\,r\cos\theta=p$ with $x=r\cos\theta$, $y=r\sin\theta$: then
$r=p-ex$, and squaring, $x^2+y^2=(p-ex)^2$, that is
\[
(1-e^2)\,x^2+2pe\,x+y^2=p^2 .
\]
For $e<1$ both squares have positive coefficients: an \emph{ellipse}. Completing the square shows that
its centre lies at $x=-pe/(1-e^2)$ and its semi-major axis is $a=p/(1-e^2)$, so the origin sits a distance
$ea$ from the centre, which is exactly where a focus of an ellipse lies: the Sun is at a focus. For
$e=1$ the $x^2$ term vanishes and the curve is a parabola; for $e>1$ it is a hyperbola.
\stepn{6} \textbf{Connect with the energy.} At perihelion ($\theta=0$) the speed is $L/(mr)$ with
$r=p/(1+e)$; substituting into $E=\tfrac12mv^2-GMm/r$ gives
\[
E=\frac{G^2M^2m^3}{2L^2}\,(e^2-1).
\]
So $E<0\Leftrightarrow e<1$ (ellipse), $E=0\Leftrightarrow e=1$ (parabola), $E>0\Leftrightarrow e>1$
(hyperbola): precisely the table of \S\ref{sec:escape}, now proved.
\end{derivbox}

\begin{remarkbox}{What the proof shows on the way}
Step~3 says something striking on its own: as the planet goes round, the tip of its velocity vector,
drawn from a fixed point, traces a \emph{circle} of radius $GMm/L$, centred at $\vc{C}$. Hamilton named
the curve traced by a velocity vector the \emph{hodograph}. Step~2 used the inverse-square form; for any
other power the $r^2$ would not cancel and the method would fail. For a force that departs slightly from
$1/r^2$, a bound orbit does not quite close: it is an ellipse that slowly turns. That is why the part of
Mercury's turning ellipse that the other planets could not explain (\S\ref{sec:beyondnewton}) was such
strong evidence that something beyond Newton's $1/r^2$ was at work.
\end{remarkbox}

\section{Summary}\label{sec:summary17}
Kepler's three laws are consequences of Newton's law of gravitation. The law of areas is the conservation of
angular momentum under a central force. The law of periods follows from equating gravity to the centripetal
force. The law of ellipses follows from the inverse-square form of the force (\S\ref{sec:ellipse}), and
the sign of the energy decides whether the conic is an ellipse, a parabola or a hyperbola.

\begin{itemize}
\item Equal areas in equal times $\Leftrightarrow$ $L=mvr_\perp$ constant, so a planet moves fastest at
      perihelion.
\item Circular orbit: $v_o=\sqrt{GM/r}$, $\;T=2\pi\sqrt{r^3/GM}$, so $T^2\propto r^3$ (and
      $T^2\propto a^3$ for an ellipse).
\item Total energy of a circular orbit $E=-\dfrac{GMm}{2r}=-K$; of an elliptical orbit,
      $E=-\dfrac{GMm}{2a}$, with $a=\tfrac12(r_p+r_a)$.
\item At perihelion and aphelion $v_pr_p=v_ar_a$; an ellipse has $r_p=a(1-e)$, $r_a=a(1+e)$.
\item Escape speed is $\sqrt2$ times the circular orbital speed at the same radius.
\item A geostationary satellite has a period of one sidereal day and orbits at about \SI{35800}{km} above the
      equator.
\end{itemize}

\section{Problem Types}\label{sec:types17}
Six types cover orbit questions. Nearly all of them follow from setting gravity equal to the centripetal
force, which gives $v_o$, $T^2\propto r^3$ and $E=-GMm/2r$, together with angular-momentum conservation for
orbits that are not circular. The most efficient approach is usually to work with ratios, so that $G$ and
$M$ cancel and no constants need be looked up.

\archq{Orbit problems: gravity supplies the centripetal force}
{``Find the speed, period or radius of a satellite at a given height.''}
{Set gravity $=$ centripetal ($GMm/r^2=mv^2/r$), with $r=R+h$ measured from the \emph{centre}, and use
$GM=gR^2$ to eliminate $G$ and $M$ together.}
{Using the height $h$ where the formulas need the radius $R+h$.}
{A satellite orbits at a height $R$ above the Earth's surface ($R$ is the Earth's radius, $g$ the surface
gravity). Its period is:}
{$2\pi\sqrt{R/g}$}{$4\pi\sqrt{R/g}$}{$8\pi\sqrt{R/g}$}{$4\pi\sqrt{2R/g}$\archans{D}{The orbital
radius is $r=R+h=2R$, so $T=2\pi\sqrt{r^3/GM}=2\pi\sqrt{8R^3/gR^2}=4\pi\sqrt{2R/g}$, about
\SI{4.0}{h}. Option (A) takes $r=h=R$ (height instead of radius), which is the close-orbit period;
(B) scales that period in proportion to $r$; (C) scales it as $r^2$. Only $T\propto r^{3/2}$, a factor
$2\sqrt2$, is right.}}

\archq{Kepler's third law as a ratio}
{``A satellite's orbital radius is changed by a factor; what happens to its period?''}
{$T^2\propto r^3$, so $\dfrac{T_2}{T_1}=\left(\dfrac{r_2}{r_1}\right)^{3/2}$. Never compute $GM$; take
the ratio. The same law relates any two planets around the same star.}
{Using $T\propto r$ or $T\propto r^2$. The exponent is $3/2$, and questions are designed so that the
wrong exponent gives one of the distractors.}
{If the orbital radius of a satellite is doubled, its period becomes:}
{$2\sqrt2\,T$}{$2T$}{$4T$}{$8T$\archans{A}{$T\propto r^{3/2}$, so doubling $r$ multiplies $T$ by
$2^{3/2}=2\sqrt2\approx2.83$. Options (B), (C) and (D) come from the wrong exponents $1$, $2$ and $3$.}}

\archq{Orbital speed and its dependence on radius}
{``Find the speed of a satellite in a given orbit, or compare two orbits.''}
{From the first type, $\boxed{v=\sqrt{GM/r}}$, that is, $v\propto1/\sqrt r$. A \emph{higher} orbit is a \emph{slower}
orbit, and for a close orbit $v=\sqrt{gR}\approx\SI{7.9}{km/s}$.}
{Thinking a higher satellite moves faster because it has more energy. It has more \emph{total} energy but
\emph{less} speed: the potential energy rises by twice the energy added, and the kinetic energy falls.}
{A satellite is moved to an orbit of four times the radius. Its orbital speed becomes:}
{$4v$}{$2v$}{$v/2$}{$v/4$\archans{C}{$v\propto1/\sqrt r$, so quadrupling $r$ halves $v$.}}

\archq{Energy of a satellite in orbit}
{``Find the kinetic, potential and total energy of an orbiting satellite, or the energy to raise it.''}
{For a circular orbit the three are locked together:
$\boxed{K=+\dfrac{GMm}{2r},\quad U=-\dfrac{GMm}{r},\quad E=-\dfrac{GMm}{2r}}$
so $U=2E$ and $K=-E$. The energy to move between orbits is the \emph{difference} of the $E$ values.}
{Forgetting the minus signs. A bound orbit always has $E<0$; $E=0$ is exactly the escape condition.}
{For a satellite in a circular orbit, the relation between kinetic energy $K$ and total energy $E$ is:}
{$K=E$}{$K=2E$}{$K=-2E$}{$K=-E$\archans{D}{$K=+\dfrac{GMm}{2r}$ and $E=-\dfrac{GMm}{2r}$, so $K=-E$.
Equivalently, the total energy is minus the kinetic energy.}}

\archq{Geostationary satellite}
{``What are the conditions for a satellite to appear fixed above a point on Earth?''}
{Four conditions together: period of one sidereal day (\SI{23}{h}~\SI{56}{min}, usually rounded to
\SI{24}{h} in problems), orbit in the \textbf{equatorial} plane, motion \emph{west to east} (same sense as the Earth's spin), and circular. These force a unique
height of about \SI{36000}{km}, found from $T^2=\dfrac{4\pi^2r^3}{GM}$.}
{Thinking any \SI{24}{h} orbit is geostationary. An inclined \SI{24}{h} orbit traces a figure-of-eight in
the sky, not a fixed point.}
{A geostationary satellite must orbit:}
{over the poles}{in the equatorial plane, west to east}{in any plane with $T=\SI{24}{h}$}{at any height
above the equator\archans{B}{Only an equatorial, west-to-east, \SI{24}{h} circular orbit keeps the
satellite above a fixed ground point, which fixes its height at roughly \SI{36000}{km}. Option (D)
fails because only one radius gives a one-day period.}}

\archq{Kepler's second law as angular momentum conservation}
{``Compare a planet's speed at perihelion and aphelion,'' or ``show equal areas are swept in equal
times.''}
{Gravity is a \emph{central} force, acting along $\vr$, so $\vtau=\vr\ex\vF=\vc0$ and angular momentum is
conserved. Hence $mvr_\perp=$ constant, and at the two apses, where $\vv\perp\vr$,
$\boxed{v_1r_1=v_2r_2}$: closest means \emph{fastest}.}
{Trying to use energy conservation alone. It is true but insufficient: the clean ratio comes from
angular momentum.}
{A planet's distance from the Sun at perihelion is one third that at aphelion. The ratio of its speeds
(perihelion: aphelion) is:}
{$1:3$}{$3:1$}{$1:9$}{$9:1$\archans{B}{$v_1r_1=v_2r_2$, so $\dfrac{v_p}{v_a}=\dfrac{r_a}{r_p}=3$. The
planet moves three times faster when three times closer.}}

\section{Challenge Problems}\label{sec:challenge17}
The three problems below draw on \chsref{ch:gravitation}{ch:kepler}. The first shows how much can be done with Kepler's
third law as a pure ratio; the second and third turn on the energy accounting of orbits and escape, where the
sign of the potential energy carries the physics.

\begin{challengebox}{Kepler as a ratio, no constants (builds on Kepler's third law)}
Two satellites orbit Earth: a low one at $r_1=\SI{7000}{km}$ and a geostationary one at
$r_2=\SI{42200}{km}$. Without using $G$, $M$, or $g$, find how many times longer the geostationary
period is, and how the speeds compare.
\challengeans{Kepler's third law as a pure ratio needs no constants:
\[
\frac{T_2}{T_1}=\left(\frac{r_2}{r_1}\right)^{3/2}=\left(\frac{42200}{7000}\right)^{3/2}=(6.03)^{1.5}\approx14.8.
\]
So the geostationary satellite takes about $15$ times as long per orbit. For speeds,
$v=\sqrt{GM/r}\propto1/\sqrt r$, so $\dfrac{v_1}{v_2}=\sqrt{\dfrac{r_2}{r_1}}=\sqrt{6.03}\approx2.5$: the
low satellite moves two-and-a-half times faster.}
\end{challengebox}

\begin{challengebox}{Escaping from Mars (builds on escape velocity)}
Mars has mass $0.107 M_E$ and radius $0.532 R_E$. Using Earth's escape velocity of \SI{11.2}{km/s}, find
Mars's escape velocity without looking up $G$ or any masses.
\challengeans{Since $v_e=\sqrt{2GM/R}$, form the ratio so $G$ cancels:
\[
\frac{v_{e,\text{Mars}}}{v_{e,\text{Earth}}}=\sqrt{\frac{M_{\text{Mars}}/R_{\text{Mars}}}{M_E/R_E}}
=\sqrt{\frac{0.107}{0.532}}=\sqrt{0.201}=0.448.
\]
So $v_{e,\text{Mars}}=0.448\times11.2\approx\SI{5.0}{km/s}$, less than half Earth's. The low escape
speed is a large part of why Mars has only a thin atmosphere: over billions of years the lighter gas
molecules in the fast tail of the speed distribution have exceeded \SI{5}{km/s} and leaked away (the
solar wind, unhindered by a global magnetic field, has stripped away more).}
\end{challengebox}

\begin{challengebox}{The energy to raise a satellite (builds on orbital energy)}
A satellite of mass $m$ is in a circular orbit at $r=2R$ (twice the Earth's radius). Find its orbital
speed in terms of surface $g$ and $R$, and the total energy that must be supplied to place it in that
orbit from rest on the surface (ignore the Earth's rotation and air resistance). Take
$g=\SI{9.8}{m/s^2}$ and $R=\SI{6.4e6}{m}$ for the numerical value of the speed.
\challengeans{At $r=2R$, orbital speed $v_o=\sqrt{\dfrac{GM}{2R}}=\sqrt{\dfrac{gR}{2}}$ (using $GM=gR^2$),
about \SI{5.6}{km/s}. The total orbital energy is
$E_{\text{orbit}}=-\dfrac{GMm}{2(2R)}=-\dfrac{GMm}{4R}=-\dfrac{mgR}{4}$. On the surface, at rest, the
energy is purely potential, $E_{\text{surface}}=-\dfrac{GMm}{R}=-mgR$. The energy that must be supplied
is
\[
\Delta E=E_{\text{orbit}}-E_{\text{surface}}=-\frac{mgR}{4}-(-mgR)=\frac{3}{4}mgR.
\]
Of this, $\tfrac12mgR$ raises the body against gravity (the work found in \S\ref{sec:gpe}) and the
remaining $\tfrac14mgR$ is its orbital kinetic energy.}
\end{challengebox}

\section*{Looking Back on the Volume}
\addcontentsline{toc}{section}{Looking Back on the Volume}

\textbf{Part E turned the whole volume towards the sky.}
\begin{itemize}
\item \textbf{Gravitation (\chref{ch:gravitation}).} One inverse-square, attractive, central force, with one
constant $G$. The shell theorem, proved by adding scalar potentials, lets a sphere act as a
point at its centre, so in a uniform planet $g$ peaks at the surface and falls both ways. Potential energy
$U=-GMm/r$ is negative because the pair is \emph{bound}, $mgh$ is its small-height limit, and the sign of
$E=K+U$ decides an orbit's fate, with escape velocity $\sqrt{2gR}$ the speed at which $E$ reaches zero.
Newton's law is itself the weak-field, low-speed limit of general relativity.
\item \textbf{Orbits (\chref{ch:kepler}).} Kepler's second law \emph{is} angular-momentum conservation
(\chref{ch:rotation}) applied to a central force; his third, $T^2\propto r^3$, falls out of setting gravity
equal to the centripetal requirement (\chref{ch:circdyn}), with the exponent $3=1+2$ a fingerprint of the
inverse-square law; and his first follows from the same force by Hamilton's argument. Angular momentum and
energy together give the speeds on an ellipse, whose energy $-GMm/2a$ depends only on its size. A
satellite is a projectile (\chref{ch:projectiles}) that keeps missing the Earth, weightless because it is
in free fall (\chref{ch:newton}), not because gravity is gone.
\end{itemize}

\medskip
\textbf{The volume as a whole.} The five parts form one argument. Part~A built the language: the
mathematical toolkit of calculus and vectors, measurement with its dimensions and uncertainties, and the
kinematics of a line, a projectile, a moving frame and a circle, all of which describe motion without
asking its cause. Part~B supplied the cause, $\vF=d\vp/dt$ applied through a free-body diagram, and with
it solved pulleys and springs, friction and circular motion, adding pseudo-forces for accelerating frames.
Part~C found the shortcuts: energy and momentum, conserved quantities that connect the start of a motion
to its end without the middle, with the coefficient of restitution measuring how much kinetic energy a
collision keeps. Part~D gave bodies size and spin, showed that a system's centre of mass moves as a single particle
under the external force alone, and built the rotational twin of every translational law, with angular momentum as the third
conserved quantity. Part~E then needed only one new ingredient, the law of gravitation, to explain a
falling apple, the Moon and the planets with those same tools. A sliding block, a spinning skater and a
planet's orbit are, in the end, one subject: Newton's laws, written in components, and read through the
quantities they conserve.
